\section{Fixed-precision approximation and accepting computations}
\label{sec:conclusion}

We prove Theorem~\ref{thm:fixed-precision}: for each fixed machine
$\mathcal M$ and integer $d\ge1$, a deterministic logarithmic-space
transducer approximates $p_{\mathcal M}(x)$ within $(n+2)^{-d}$ in
polynomial time, with logarithmic output length. The hierarchy gives
a deterministic reward close to the acceptance probability. The
estimator's moment bound gives a majority of accurate environments,
and total evaluation on every environment permits median enumeration.
Section~\ref{sec:deduction} already derives the class equality, promise
separation, and unary-precision approximation from this result.

\subsection{Choosing constants and finite ranges}
\label{subsec:final-constants}

The construction has the following order of choices.
\begin{enumerate}
\item Fix $H,h_0,s=8$ and the fixed number of independent sampling matrices.
The finite-action gap is uniform in every later modulus and bit length.
\item Choose the correction thresholds and cutoff, then the iteration
count $k_*$ from the factorial bound and its resulting uniform
magnitude constants. These estimates hold for every column multiplier
in $[0,1]$.
\item Choose $D_0$ for the entry bound reduction, the routing thresholds,
and the finite ladders of entry and cutoff thresholds. The row-cutoff
rank $K$ can then be chosen from the row-mass bounds alone.
\item Choose the detector rank $J$ large enough for the incidence
bound on extra copies. Choose $R_*$ large enough for overload
detection, and finally the incoming rank $K_c$ large enough for
correction normalization. The earlier row-cutoff choice does not depend
on $K_c$.
\item Fix the bounded-output constant $C_b$, the scale cutoff $c_0$ and
precision coefficient $K_{\rm prec}$, and the prefix mean exponent $B_1$.
Choose the buffer floor and capacity exponents $c_1,c_2$ and the noise
decay $K_n$ so that all buffer sums converge. Choose the rank coefficient
$A_*$ of Equation~\eqref{eq:est-precision} and the walk-length
coefficient $\ell$ sufficiently large for the small constants in the
comparison estimation induction.
\item Fix the crude range and propagation constants, then the block
gain coefficients $D_2,D_1$ and a sufficiently small rational block
width coefficient $\epsilon$. These choices close the coupled
fingerprint induction of Theorem~\ref{thm:fp-estimators}.
\item Choose the fixed operation-stratum and allowance coefficients
and the guard constants in Theorem~\ref{thm:ctrl-realization}.
Their finite enlargement does not alter the analytical tables or
the distribution of the environment.
\item Choose the linear top-budget coefficient, the field-size
exponent, and the minimum padding of $B$.
\end{enumerate}
Within each step only finitely many strict inequalities between
fixed constants are needed. The proofs in the corresponding sections
give these inequalities. In particular, the row rank is fixed before
the incoming rank, the block constants follow the crude error
constants, and the field is chosen after the reachable resource
ranges. There is no dependence of an earlier averaging walk gap on that
last choice.

The coefficient $\epsilon$ in this order is the fixed fingerprint
block-width coefficient. Once it has been chosen, increasing the fixed
linear top budget increases the number of seed blocks but leaves the
averaging walk's gap uniform: all earlier seed bits form one free component in
Lemma~\ref{lem:env-gap}. The storage estimate
\eqref{eq:fp-salt-space} therefore remains $O(B)$.

The allowed family in Section~\ref{subsec:est-tasks} bounds the stage
by $L$, each inherited suffix length by $L-l$ at stage $l$, and every
reciprocal-rate logarithm by
\begin{equation}\label{eq:final-admissible}
HL+C_{\rm rate}(M_{\rm top}-M).
\end{equation}
The closure check after Lemma~\ref{lem:est-paths} verifies both this
finite range and the task-specific inequalities $u\le f$,
$a\le K_f f$ and $t\le a$, including stage resets, preparations,
detectors and reward calls.  The finitely many local operations keep
their task tuple and use the stated widths and fixed subranks.
The controller guards allow these designated executions.

There is consequently a fixed $C_P$ such that a prime
\[
2^{C_PB}\le P\le 2^{C_PB+1}
\]
is larger than all identifiers, makes the radix invertible, and gives
$P2^{-k}\ge1$ for every rate ever requested. Bertrand's postulate
provides such a prime, as it provides the prime of
Lemma~\ref{lem:fp-hash}.
Testing consecutive integers by trial division finds the prime
using $O(B)$ space. The search range and running time are
$2^{O(B)}$. Field arithmetic and fixed-dimensional rank tests also
use $O(B)$ space.

\subsection{A dyadic approximation at the specified start}

\begin{lemma}[Conditioning on one start]\label{lem:final-majority}
Let $V$ be finite, let $W\in[0,1]^V$, and use the sampling matrix
indicators $X_f$ at an admissible rate $f$, with marginal $\pi_f>0$.
Suppose an estimator satisfies
\[
\frac1N\mathbb E\sum_{v\in V}\frac{X_f(v)}f
 |\bar W(v)-W(v)|^8\le C_b^8 2^{-8p}.
\]
For any specified $z\in V$ and every $\eta>0$,
\begin{align}
\mathbb E\bigl[|\bar W(z)-W(z)|^8\mid X_f(z)=1\bigr]
 &\le N\frac f{\pi_f}C_b^8 2^{-8p},
 \label{eq:final-conditional}\\
\Pr\bigl[|\bar W(z)-W(z)|>\eta\mid X_f(z)=1\bigr]
 &\le N\frac f{\pi_f}C_b^8 2^{-8p}\eta^{-8}.
 \label{eq:final-tail}
\end{align}
\end{lemma}

\begin{proof}
Keep only the nonnegative summand at $z$. Since
$\Pr[X_f(z)=1]=\pi_f>0$,
\[
 \frac{\pi_f}{f}
 \mathbb E\bigl[|\bar W(z)-W(z)|^8\mid X_f(z)=1\bigr]
 =\mathbb E\!\left[\frac{X_f(z)}f|\bar W(z)-W(z)|^8\right]
 \le N C_b^8 2^{-8p}.
\]
Dividing by $\pi_f/f$ proves
Equation~\eqref{eq:final-conditional}. In particular, choosing one
specified coordinate costs the factor $N$ from the integrated norm.
Markov's inequality gives Equation~\eqref{eq:final-tail}.
\end{proof}

\begin{proof}[Proof of Theorem~\ref{thm:fixed-precision}]
Apply Lemma~\ref{lem:model-reduction} to $\mathcal M$. Set
\[
 r=d\lceil\log_2(n+2)\rceil,\qquad \Delta=2^{-r},\qquad
 L=\left\lceil\frac{C_0B+r+3}{H-1}\right\rceil,
\]
where $|V_0|\le2^{C_0B}$ is an explicit configuration bound.
Since $d$ is fixed, $r,L=O(B)$. Let $z$ be the copy with every appended digit zero of
the start at stage $L$. Proposition~\ref{prop:hier-reward-target}
gives
\begin{equation}\label{eq:final-hierarchy-accuracy}
 0\le p_{\mathcal M}(x)-W_L(z)\le\Delta/8.
\end{equation}
This bound holds for every acceptance probability. Choose a fixed
$C_N$ large enough to bound all the resulting stage domains by
$N=2^{C_NB}$.

For the final reward task take $a=u=f=b_L$. Its rank satisfies
$i_*=O(L+1)$, and Equation~\eqref{eq:est-precision} reads
\[
 p=M_{\rm top}/K_{\rm prec}-HL-A_*i_*.
\]
By Lemma~\ref{lem:env-ranks}, fix an integer $C_{\rm tail}\ge1$ bounding
$(f/\pi_f)C_b^8$ throughout the admissible range. Choose the fixed
linear coefficient of $M_{\rm top}=O(B)$ large enough that
\begin{equation}\label{eq:final-precision-choice}
 p\ge r+3+\frac{C_NB+\lceil\log_2 C_{\rm tail}\rceil+3}{8}.
\end{equation}
The top-budget and field-size choices in the preceding subsection
allow this enlargement; Equation~\eqref{eq:final-admissible} still
bounds every descendant rate. Theorem~\ref{thm:fp-estimators}
supplies the reward moment estimate. Lemma~\ref{lem:final-majority},
with $\eta=\Delta/8$, now shows that
\begin{equation}\label{eq:final-good-environments}
 \Pr\bigl[|\bar W_L(z)-W_L(z)|>\Delta/8
       \mid X_f(z)=1\bigr]\le1/8.
\end{equation}
Thus the factor $N$ lost on selecting one coordinate has been paid
by the statistical precision. Explicitly, the bound in
Equation~\eqref{eq:final-tail} is at most
\[
 C_{\rm tail}\,2^{C_NB-8p+8r+24}\le 2^{-3},
\]
by Equation~\eqref{eq:final-precision-choice}.

We next turn the estimator values into a finite list of dyadic numbers.
For each master environment, Theorem~\ref{thm:arith-interface}
gives the digits $w_j\in\{-6,\ldots,6\}$ of its one fixed value
$\bar W_L(z)$. Structural choices do not depend on the requested
digit. Put $K=r+6$, compute
\[
 A_\sigma=\sum_{j=0}^{K}w_j2^{K-j},\qquad
 a_\sigma=\min\{2^K,\max\{0,A_\sigma\}\},\qquad
 v_\sigma=a_\sigma2^{-K}.
\]
The integer $A_\sigma$ is obtained by starting at zero and applying
$A\leftarrow2A+w_j$ in order, so it uses $O(K)$ bits.
By Equation~\eqref{eq:arith-tail}, the truncation error is at most
$6\,2^{-K}<\Delta/8$. Clipping cannot increase error relative to
$W_L(z)\in[0,1]$. Equations~\eqref{eq:final-hierarchy-accuracy}
and~\eqref{eq:final-good-environments} imply that, on more than
three quarters of the environments conditioned on $X_f(z)=1$,
\begin{equation}\label{eq:final-dyadic-accuracy}
 |v_\sigma-p_{\mathcal M}(x)|<3\Delta/8.
\end{equation}

Enumerate the conditional law exactly. Enumerate each field-matrix
entry in $\{0,\ldots,P-1\}$ and each raw seed bit string, retaining
precisely those encodings whose positive-degree columns are independent
and for which $X_f(z)=1$. The retained encodings are equally weighted
and give the law in Equation~\eqref{eq:final-good-environments}.
In particular, enumerate the original seed strings, not their reduced
fingerprint parameters. Different raw strings that induce the same
fingerprint remain separate encodings; their multiplicities reproduce
the original probability law. At least one encoding is retained: an independent column tuple
exists, and its independent constant column can place the designated
evaluation in the nonempty retention box. The matrices have fixed
dimensions and $\log P=O(B)$; together with
Equation~\eqref{eq:fp-salt-space}, this gives $O(B)$ bits per master
environment and at most $2^{O(B)}$ encodings to enumerate.

Let $E\ge1$ be the number of retained encodings, found by one pass,
and let $h=\lfloor E/2\rfloor+1$. Binary-search the integer interval
$[0,2^K]$ for the least $a$ such that at least $h$ retained encodings
satisfy $a_\sigma\le a$. Each count is computed by a fresh pass,
recomputing each numerator $a_\sigma$. This finds the upper median
of the finite multiset, including all repetitions. Since more than
three quarters of its members satisfy
Equation~\eqref{eq:final-dyadic-accuracy}, so does its median: fewer than
$E/4$ values can lie below the indicated interval, and fewer than
$E/4$ can lie above it. The algorithm retains only the current count
and search registers; it never stores the multiset.
Output the integers $(a,K)$ representing $z_x=a2^{-K}$. Then
\[
 |z_x-p_{\mathcal M}(x)|<3\Delta/8
       \le(n+2)^{-d}.
\]

We account for the memory held during this computation. The master
environment, current cursor, catalyst, shared scratch, and suspended
controller records together use $O(B)$ bits by
Theorem~\ref{thm:ctrl-realization}, because $M_{\rm top},K=O(B)$.
The outer computation keeps only a fixed number of additional
$O(B)$-bit registers: its top parameters and designated start,
the environment enumeration index, $E$, $h$, the current count,
the two binary-search endpoints and candidate, the digit index,
and the signed numerator accumulator. Each count is bounded by the
number of environment encodings, so it also has $O(B)$ bits.
These outer registers remain live while a digit query runs; in
particular, both its environment-enumeration state and its numerator
accumulator are included. They are present once for the whole
computation, not once per recursive level. Every query terminates on every environment,
including the inaccurate environments, by the same realization
theorem. The output is written only after the median is computed.

Each digit query is a halting deterministic computation with $O(B)$
work bits and physical input heads confined between the real endmarkers.
Counting its finite control, work-tape contents, and all work and input
head positions gives at most $2^{O(B)}$ complete configurations on its
fixed input. A repeated complete configuration would
force a deterministic loop, so totality bounds the query by
$2^{O(B)}$ steps. There are
$2^{O(B)}$ environment encodings, $O(K+1)$ binary-search passes,
and $K+1$ digit queries per numerator. Their product is still
$2^{O(B)}=(n+2)^{O(1)}$. All finite parameters are generated
uniformly, and all constants are fixed for $(\mathcal M,d)$.
\end{proof}

\subsection{Certified accepting computations}

We return to fixed inverse-polynomial accuracy to construct an accepting
computation when the machine's acceptance probability has an
inverse-polynomial lower bound.
For a transducer, the conclusion concerns outputs certified by an
accepting run: its output tape must be write-only, so past emitted
symbols do not influence later transitions.

\begin{corollary}[Certified accepting computations]
\label{cor:certified-search}
Fix a randomized logarithmic-space machine $\mathcal M$ that halts
within a fixed polynomial bound $T(n)\ge1$ on every input and every
coin sequence, and fix an integer $c\ge0$. There is a uniform
deterministic logarithmic-space, polynomial-time transducer that
outputs an accepting computation of $\mathcal M$ whenever
\[
 p_{\mathcal M}(x)\ge(n+2)^{-c}.
\]
On every other input it outputs either an accepting computation or
a failure symbol. If every accepting run of $\mathcal M$ emits a
valid witness for a specified relation, the transducer can instead
output such a witness with the same guarantee.
\end{corollary}

\begin{proof}
Give halted configurations constant self-transitions through time
$T(n)$. A configuration and its elapsed time use $O(\log(n+2))$
bits. There is a fixed randomized logarithmic-space continuation
machine which, on a valid encoding of $(x,v,t)$, starts from
configuration $v$ and simulates the remaining $T(n)-t$ steps with
fresh coins. It checks the fixed work-space box and time bound,
rejecting malformed encodings and any attempted violation. Its
acceptance probability on the reachable configurations is the
continuation value $p_t(v)$, and it has a
polynomial worst-case time bound in its input length.

Write $\gamma=(n+2)^{-c}$. Choose a fixed exponent $d$ large enough
that Theorem~\ref{thm:fixed-precision}, applied to the continuation
machine, approximates every queried value within
$\delta\le\gamma/(8T(n))$. For an explicit choice, fix $C_T\ge1$ and an integer $a_T\ge0$
such that $T(n)\le C_T(n+2)^{a_T}$ for every $n$, and take
\[
 d\ge a_T+c+3+\lceil\log_2 C_T\rceil.
\]
If $m$ is the continuation input length, then $m\ge n$, and
\[
 (m+2)^{-d}\le(n+2)^{-d}
 \le\frac{(n+2)^{-c}}{8C_T(n+2)^{a_T}}
 \le\frac\gamma{8T(n)}.
\]
The extra fields have logarithmic length. We supply them virtually
from saved configuration and time registers, so the approximation
and those registers together use logarithmic space.

Starting at the initial configuration, at each time estimate the
continuation value of every positive-probability successor and
choose one with maximum estimated value, breaking ties by a fixed
port order. Since $p_t(v)$ is the weighted mean of the successor
values, the chosen successor $v'$ satisfies
\[
 p_{t+1}(v')\ge\max_w p_{t+1}(w)-2\delta
                  \ge p_t(v)-2\delta.
\]
On an input satisfying the stated lower bound, after $T(n)$ choices
the terminal continuation value is at least
$\gamma-2T(n)\delta\ge3\gamma/4>0$. A terminal continuation value
is zero or one, so this computation accepts.

First run this construction without writing output. If its terminal
state rejects, output the failure symbol. Otherwise repeat the same
deterministic choices and stream the configurations through the original
first halting state, or the symbols emitted by the simulated transducer.
The artificial self-transitions after halting are omitted from the output.
This second pass follows the same accepting computation. Consequently no rejecting computation
is output, even outside the probability promise.

At any time the current, candidate, and best successor configurations,
the time counter, a constant number of dyadic scores, the virtual
input address, and the approximation workspace occupy a constant
number of $O(\log(n+2))$-bit records. The finite successor lists
have constant size, and there are only $T(n)$ steps in each pass,
so the total running time is polynomial. Each approximation has a
deterministic guarantee at its queried configuration; no simultaneous
accuracy event for all adaptive queries is needed.
\end{proof}
